\begin{ExerciseBlock}
\ExerciseHeading{06}{Gleichförmige Kreisbewegung}{multidimensional_kinematics_uniform_circular_motion}
Eine Spielzeugeisenbahn fährt auf einer kreisförmigen Schienenstrecke mit Radius $R = 2 \,$m und einer Winkelgeschwindigkeit $\omega = 2,5 \, s^{-1}$.

\begin{enumerate}
	\item Berechnen Sie die Umlaufdauer des Zuges $T$.

	\item Mit welcher Frequenz $f$ (Drehzahl) vollziehen sich die Umläufe?

	\item Berechnen Sie die Fahrtgeschwindigkeit des Zuges.

	\item Wie groß ist die Zentripetalbeschleunigung $a_{zp}$, die der Zug durch die Schienen erfährt?

	\item Angenommen der Zug fährt von oben gesehen im Uhrzeigersinn, welche Richtung hat die Winkelgeschwindigkeit $\vec{\omega}$?

\end{enumerate}
\end{ExerciseBlock}
\UseExerciseSolution{multidimensional_kinematics_uniform_circular_motion}
